\section{Relation to Algorithmic Regulation}\label{sec:ruffini}

\citet{Ruffini2026} reformulates the good regulator principle in algorithmic
terms: a regulator that materially simplifies the realized world trajectory is
unlikely to be independent of that world. In Ruffini's setting, regulation is
read through compression of the regulated output. The question here is
different: not how much compression a regulator achieves on a realized episode,
but how long a compressed internal model remains valid when the world itself
drifts.

Tracking under drift can be interpreted
as \emph{online compression of a non-stationary process}. A finite-memory model
compresses recent observations into a compact state, but that compression ages.
The cube-root law in \cref{prop:coherence_delay} quantifies the resulting
trade-off between statistical efficiency and temporal validity.

The decomposition $(\Pot,\Act,\Conv)$ then serves as an explicit model of the
three channels that matter in closed-loop regulation: what the system currently
represents, what it does, and how it updates. Ruffini studies the compression
value of regulation on individual episodes, whereas the analysis here tracks
the stability of that compression under drift and feedback.

\subsection{Coercive Masking}\label{sec:coercive}

The most important point of contact with algorithmic regulation is the masking
regime. A controller can keep the realized trajectory simple either because it
tracks the world well or because it pushes the world toward a stale model. The
latter mechanism is invisible to a diagnostic that looks only at model-world
fit.

\begin{definition}[Coercive Regime]\label{def:coercive}
An information system $\IS = (\Pot, \Act, \Conv)$ operates in the
\emph{coercive regime} if:
\begin{enumerate}[label=\textnormal{(\roman*)}]
  \item $\Act_t$ is action-effective (\cref{def:actioneffective}) with
    coupling coefficient $\alpha > 0$: the transition kernel is
    $P(s_{t+1} \mid s_t, o) = \mathcal{N}(\mu^*(t) + \alpha(o - \mu^*(t)),
    \sigma_w^2)$;
  \item $\Pot$ is static (\cref{prop:failP}): $\hat\mu_t \equiv \hat\mu_0$
    and $\Act_t(\Pot) \equiv \hat\mu_0$;
  \item the exogenous drift is bounded:
    $|\zeta| < \alpha \cdot |\hat\mu_0 - \mu_0|$.
\end{enumerate}
\end{definition}

In the coercive regime, the system's stale actions continuously pull
$\mu^*$ toward the frozen estimate $\hat\mu_0$:
$\mu^*(t+1) = \mu^*(t) + \zeta + \alpha(\hat\mu_0 - \mu^*(t))$,
which is a mean-reverting process with equilibrium $\hat\mu_0 + \zeta/\alpha$.
When $|\zeta/\alpha|$ is small, $\mu^*$ is held near $\hat\mu_0$,
suppressing $V_P$ and making the system appear healthy under the standard score.

\begin{corollary}[Coercive Masking]\label{cor:coercive}
In the coercive regime (\cref{def:coercive}), the standard score
$\sigma_P(t) = 1/(1 + V_P(t))$ can fail to detect FM-1:
$\sigma_P(t) \approx 1$ for all $t \leq t_{\mathrm{mask}}$, where
$t_{\mathrm{mask}} = O(\alpha^{-1} \log(\alpha/\zeta))$ is the
\emph{coercive horizon}. Beyond $t_{\mathrm{mask}}$, exogenous drift overwhelms
the coercive capacity and $V_P$ diverges.
\end{corollary}

\begin{proof}[Proof sketch]
Under condition (iii) of \cref{def:coercive}, the gap
$\delta_t = \mu^*(t) - \hat\mu_0$ satisfies
$\delta_{t+1} = (1 - \alpha)\delta_t + \zeta$. The fixed point is
$\delta^* = \zeta / \alpha$, so
$V_P(t) \leq \tfrac{1}{2}(\zeta/\alpha)^2 + \sigma_w^2/(2\alpha)$ in
expectation, masking FM-1 while the coercive condition holds. Once the drift
accumulates beyond the available control authority, the frozen model can no
longer hold the environment in place and the underlying mismatch reappears.
\end{proof}

\begin{remark}[Compression by adaptation vs. compression by forcing]
From the algorithmic-regulation viewpoint, coercive masking matters because it
can improve realized compressibility without improving the model. The external
trajectory becomes simpler, but for the wrong reason: the controller is forcing
the world into alignment rather than updating its internal state to match a
changing world.
\end{remark}

\begin{definition}[Effort-Corrected Score]\label{def:tci_effort}
The \emph{effort signal} at time $t$ is any monotone proxy for the corrective
actuation burden required to keep the world aligned with the system's present
output. A generic observable surrogate is
\[
  E_t^{\mathrm{obs}} \coloneqq |\Act_t(\Pot) - \Act_{t-1}(\Pot)|,
\]
the magnitude of change in the system's actuation. In controlled synthetic
settings, a sharper proxy is the realized coercive effort
\[
  E_t^{\mathrm{sim}} \coloneqq \alpha\,|a_t - s_t^{\mathrm{free}}|,
\]
where $s_t^{\mathrm{free}}$ is the state that would have arisen before applying
control.

The \emph{effort-corrected score} is
\[
  \ScoreE(\IS,t)
  \coloneqq
  \min\!\left(
    \sigma_P(t) \cdot e^{-\lambda E_t / E_0},\;
    \sigma_A(t),\;
    \sigma_\Phi(t)
  \right),
\]
where $E_0 > 0$ is a reference effort scale and $\lambda > 0$ is a
sensitivity parameter. $\ScoreE$ penalizes increasing actuation effort,
correcting for action-induced confounding: it discounts apparent coherence when
the environment is being held in place rather than genuinely estimated.
\end{definition}

The empirical evidence for coercive masking comes from the particle-tracker experiment
(\cref{sec:tpt_masking}). In the passive FM-1 regime, the standard score and
$\ScoreE$ coincide because no realized forcing of the world occurs.
Under moderate coupling ($\alpha = 0.3$, $\lambda = 3.0$), the same frozen
model reaches $\Score = 0.949 \pm 0.012$ while
$\ScoreE = 0.697 \pm 0.046$, showing that the apparent coherence is
being purchased by intervention.
